\section{Related work}
\label{sec:related}

\paragraph{Ultra-large virtual screening.}
Docking campaigns against make-on-demand libraries have recovered potent
ligands at hit rates that improve with library
size~\citep{lyu2019ultralarge}. The AmpC $\beta$-lactamase and D4 dopamine
receptor screens from that work are the source of two of the pools used here,
and their public release of docked scores is what makes surrogate benchmarking
on realistic distributions possible at all.

\paragraph{Surrogate models and active learning.}
MolPAL framed accelerated screening as pool-based active
learning, docking an initial batch, fitting a surrogate, and acquiring further
batches by predicted score or upper confidence
bound~\citep{graff2021molpal}. Deep Docking applies a related iterative
filtering scheme~\citep{gentile2022deepdocking}. Both report that a
surrogate trained on a small fraction of the pool recovers a large fraction of
the top-scoring molecules. Their representations are typically fingerprints or
descriptors; our question is what pretrained embeddings add relative to that
starting point.

\paragraph{Molecular representation learning.}
Pretrained molecular models differ in the modality they observe: SMILES
sequence models~\citep{ross2022molformer,soares2025smited}, graph-based
models~\citep{rong2020grover}, hypergraph-grammar
autoencoders~\citep{kajino2019mhg}, and 3D conformer
models~\citep{zhou2023unimol}. Extended-connectivity
fingerprints~\citep{rogers2010ecfp} remain a strong classical baseline. Each
carries a distinct inductive bias, and a recurring question is whether those
biases are complementary.

\paragraph{Multimodal fusion for molecular property prediction.}
\citet{martinezgalindo2026multimodality} conducted the most systematic study to
date, training over \num{1400} models across \num{7} molecular modalities,
\num{2} architectures, \num{2} fusion strategies and \num{29} benchmarks drawn
from the Therapeutics Data Commons~\citep{huang2021tdc}. They report that
additional modalities improve average performance, that late fusion
consistently outperforms intermediate fusion, and that knowledge-enhanced
pretraining yields further gains. Our work differs in three respects. First,
the target is a physics-based docking score rather than an experimental
property, so the label is cheap, abundant, and --- crucially --- has a known
dependence on molecular size that permits an explicit control. Second, the
pools are three to five orders of magnitude larger, which turns representation
dimensionality into a storage and compute decision rather than an academic one.
Third, we include a descriptor control, which that study and most of the
benchmarks it builds on do not.

\todo{Verify whether a size-controlled descriptor baseline for docking
surrogates has already been published. If so, reframe this contribution as
quantification at scale rather than as a novel control. Candidate literature:
deep-docking follow-ups, MolPAL ablations, QSAR-era work on the size dependence
of scoring functions.}
